% \usepackage[english]{babel}
\usepackage[brazil]{babel}
\usepackage[utf8]{inputenc}
\usepackage[backend=biber,style=numeric-comp,sorting=none]{biblatex}
\usepackage{comment}
\usepackage{csquotes}
\usepackage{kpfonts}
\usepackage{amsmath}
\usepackage{tcolorbox}


\usepackage{listings}
\lstset{
  language=Python,
  frame=single,  % Adds a frame around the code
  basicstyle=\ttfamily,  % Use typewriter font for code
  keywordstyle=\color{blue},  % Keywords in blue
  commentstyle=\color{gray},  % Comments in gray
  stringstyle=\color{red},  % Strings in red
  breaklines=true  % Automatically break long lines
}


%imprimir 4 por página

\usepackage{pgfpages}
\selectlanguage{portuguese}
%\pgfpagesuselayout{4 on 1}[a4paper,border shrink=5mm,landscape]
% 4 quadros

%\addbibresource{bibliografia.bib}

% 1x1 === Para os alunos === Sem pausas  % http://tex.stackexchange.com/questions/1423/is-there-a-nice-way-to-compile-a-beamer-presentation-without-the-pauses
%  \documentclass[handout]{beamer}

% 3x1 === Para os alunos === Com anotações
%  \documentclass[handout]{beamer}
%  \usepackage{handoutWithNotes}
%  \pgfpagesuselayout{3 on 1 with notes}[a4paper,border shrink=5mm]

% 8x1
%  \documentclass[handout]{beamer}
%  \usepackage{pgfpages}
%  \mode<handout>{\pgfpagesuselayout{8 on 1}[a4paper,border shrink=5mm]}%,landscape]}

% Widescreen
%  \documentclass[aspectratio = 169]{beamer}


%\usepackage{beamerthemeAmsterdam}

% Setup appearance:
\usetheme{Dresden}
\usepackage[onelanguage,ruled,lined,linesnumbered]{algorithm2e}
\SetAlFnt{\small}
\SetNoFillComment
\newcommand\mycommfont[1]{\textit{\footnotesize\ttfamily\textcolor{gray}{#1}}}
\SetCommentSty{mycommfont}
\usepackage{algorithmic}
\algsetup{linenosize=\tiny}
\SetKw{KwBy}{by}
\SetKwRepeat{Do}{do}{while}%

%portuguese

%\defbeamertemplate*{headline}{nosections}{\vskip\headheight}
\usefonttheme[onlylarge]{structurebold}
\setbeamerfont*{frametitle}{size=\normalsize,series=\bfseries}
\setbeamertemplate{navigation symbols}{}
\usecolortheme{dolphin}
\setbeamercovered{transparent} %pode ser: invisible
%\setbeamertemplate{mini frames}[tick]

%\setbeamertemplate{mini frames}[tick]

\usepackage[framemethod=tikz]{mdframed}

\definecolor{mycolor}{rgb}{0.122, 0.435, 0.698}

\newmdenv[innerlinewidth=0.5pt, roundcorner=4pt,linecolor=mycolor,innerleftmargin=6pt,
innerrightmargin=6pt,innertopmargin=6pt,innerbottommargin=6pt]{mybox}

%Para inserir rodapé
\setbeamertemplate{footline}
{
  \leavevmode%
  \hbox{%
  \begin{beamercolorbox}[wd=.3\paperwidth,ht=2.25ex,dp=1ex,center]{subsection in head/foot}%
    \usebeamerfont{author in head/foot}\inserttitle
  \end{beamercolorbox}%
  \begin{beamercolorbox}[wd=.63\paperwidth,ht=2.25ex,dp=1ex,left,leftskip=0.5em]{title in head/foot}%
    \usebeamerfont{title in head/foot}\insertsubtitle
  \end{beamercolorbox}%
  \begin{beamercolorbox}[wd=.07\paperwidth,ht=2.25ex,dp=1ex,center]{title in head/foot}%
    \usebeamerfont{title in head/foot} \insertframenumber{}/ \inserttotalframenumber
  \end{beamercolorbox}}%
  \vskip0pt%
}


% Standard packages
% \usepackage[english]{babel}
% \usepackage[latin1]{inputenc}
\usepackage{times}
\usepackage[T1]{fontenc}
\usepackage{hyperref}
\usepackage{listings}
% \usepackage{listingsutf8}
\usepackage{graphicx}
\usepackage{fancybox}
\usepackage{multirow}
\usepackage{adjustbox}
%\usepackage{xcolor,colortbl} % colorir celulas de tabelas: http://tex.stackexchange.com/questions/94799/how-do-i-color-table-columns-and-rows
%\usepackage{animate} % PNG Animado: http://www.ipgp.fr/~lucas/Contrib/animbeamer.html  /  convert Bubble_sort_animation.gif -scene 1 +adjoin Bubble_sort_animation_%0d.png

% Setup TikZ
\usepackage{lmodern}
\usepackage{xparse}
\usepackage{tikz}
\usetikzlibrary{arrows}
\tikzstyle{block}=[draw opacity=0.7,line width=1.4cm]
\usetikzlibrary{shapes.callouts} % referencias: http://tex.stackexchange.com/questions/5423/how-to-open-a-temporary-comics-like-balloon-in-a-beamer-slide
                                 %              http://tex.stackexchange.com/questions/83783/explanatory-bubbles-in-beamer
                                 %              http://tex.stackexchange.com/questions/38805/simple-speech-bubbles-arrows-or-balloon-like-shapes-in-beamer
\tikzset{
    invisible/.style={opacity=0,text opacity=0},
    visible on/.style={alt=#1{}{invisible}},
    alt/.code args={<#1>#2#3}{%
      \alt<#1>{\pgfkeysalso{#2}}{\pgfkeysalso{#3}} % \pgfkeysalso doesn't change the path
    },
}

\NewDocumentCommand{\mycallout}{r<> O{align=center, fill=cyan!20, opacity=0.8,text opacity=1} m m m m}{%
  \tikz[remember picture, overlay]
    \node[text width=#4, #2,visible on=<#1>, rounded corners,
    draw,rectangle callout,anchor=pointer,callout relative pointer={(#5)}] at (#3) {#6};
}


\newcommand{\tikzmark}[1]{\tikz[overlay,remember picture,baseline=-0.5ex] \node (#1) {};}

% Setup listings
\lstset{ %
%   inputencoding=utf8/latin1,
  %backgroundcolor=\color{white},   % choose the background color; you must add \usepackage{color} or \usepackage{xcolor}
  basicstyle=\scriptsize\ttfamily, % the size of the fonts that are used for the code
  breakatwhitespace=false,         % sets if automatic breaks should only happen at whitespace
  breaklines=true,                 % sets automatic line breaking
  captionpos=b,                    % sets the caption-position to bottom
  commentstyle=\color{gray},       % comment style
  %deletekeywords={...},            % if you want to delete keywords from the given language
  escapeinside={\%*}{*)},          % if you want to add LaTeX within your code
  %extendedchars=true,              % lets you use non-ASCII characters; for 8-bits encodings only, does not work with UTF-8
  frame=leftline,                   % adds a frame around the code
  keepspaces=true,                 % keeps spaces in text, useful for keeping indentation of code (possibly needs columns=flexible)
  keywordstyle=\color{blue},       % keyword style
  language=C++,                    % the language of the code
  %morekeywords={*,...},            % if you want to add more keywords to the set
  numbers=left,                    % where to put the line-numbers; possible values are (none, left, right)
  numbersep=6pt,                   % how far the line-numbers are from the code
  numberstyle=\tiny\color{red},    % the style that is used for the line-numbers
  stepnumber=1,                    % the step between two line-numbers. If it's 1, each line will be numbered
  rulecolor=\color{black},         % if not set, the frame-color may be changed on line-breaks within not-black text (e.g. comments (green here))
  showspaces=false,                % show spaces everywhere adding particular underscores; it overrides 'showstringspaces'
  showstringspaces=false,          % underline spaces within strings only
  showtabs=false,                  % show tabs within strings adding particular underscores
  stringstyle=\color{red},         % string literal style
  tabsize=2,                       % sets default tabsize to 2 spaces
  %title=\lstname,                  % show the filename of files included with \lstinputlisting; also try caption instead of title
  extendedchars=true,
  inputencoding=utf8,
  literate=
    {á}{{\'a}}1
    {à}{{\`a}}1
    {ã}{{\~a}}1
    {é}{{\'e}}1
    {ê}{{\^e}}1
    {í}{{\'i}}1
    {ó}{{\'o}}1
    {õ}{{\~o}}1
    {ú}{{\'u}}1
    {ü}{{\"u}}1
    {ç}{{\c{c}}}1
%   literate=%
%          {á}{{\'a}}1
%          {í}{{\'i}}1
%          {é}{{\'e}}1
%          {ý}{{\'y}}1
%          {ú}{{\'u}}1
%          {ó}{{\'o}}1
%          {ě}{{\v{e}}}1
%          {š}{{\v{s}}}1
%          {č}{{\v{c}}}1
%          {ř}{{\v{r}}}1
%          {ž}{{\v{z}}}1
%          {ď}{{\v{d}}}1
%          {ť}{{\v{t}}}1
%          {ň}{{\v{n}}}1                
%          {ů}{{\r{u}}}1
%          {Á}{{\'A}}1
%          {Í}{{\'I}}1
%          {É}{{\'E}}1
%          {Ý}{{\'Y}}1
%          {Ú}{{\'U}}1
%          {Ó}{{\'O}}1
%          {Ě}{{\v{E}}}1
%          {Š}{{\v{S}}}1
%          {Č}{{\v{C}}}1
%          {Ř}{{\v{R}}}1
%          {Ž}{{\v{Z}}}1
%          {Ď}{{\v{D}}}1
%          {Ť}{{\v{T}}}1
%          {Ň}{{\v{N}}}1                
%          {Ů}{{\r{U}}}1  
}

% Para que o codigo seja exibido com overlay, destacando partes em sequencia
% http://tex.stackexchange.com/questions/8384/how-to-make-overlay-still-work-inside-lstlisting-environment
% http://tex.stackexchange.com/questions/100522/highlight-listings-in-beamer
% http://tex.stackexchange.com/questions/45011/beamer-how-to-do-src-highlighting-when-src-is-split-with-only-command
\lstdefinestyle{highlight}{
  basicstyle=\scriptsize\ttfamily\color{black}\colorbox{yellow},
  commentstyle=\color{gray},
  keywordstyle=\color{blue},
  stringstyle=\color{red},
  backgroundcolor=\color{yellow},
}
\lstdefinestyle{overlay}{
  basicstyle=\scriptsize\ttfamily\color{black!40}, % the size of the fonts that are used for the code
  breakatwhitespace=false,         % sets if automatic breaks should only happen at whitespace
  breaklines=true,                 % sets automatic line breaking
  captionpos=b,                    % sets the caption-position to bottom
  commentstyle=\color{gray!40},       % comment style
  escapeinside={(*}{*)},          % if you want to add LaTeX within your code
  frame=leftline,                   % adds a frame around the code
  keepspaces=true,                 % keeps spaces in text, useful for keeping indentation of code (possibly needs columns=flexible)
  keywordstyle=\color{blue!40},       % keyword style
  language=C++,                    % the language of the code
  numbers=left,                    % where to put the line-numbers; possible values are (none, left, right)
  numbersep=6pt,                   % how far the line-numbers are from the code
  numberstyle=\tiny\color{red},    % the style that is used for the line-numbers
  stepnumber=1,                    % the step between two line-numbers. If it's 1, each line will be numbered
  rulecolor=\color{black},         % if not set, the frame-color may be changed on line-breaks within not-black text (e.g. comments (green here))
  showspaces=false,                % show spaces everywhere adding particular underscores; it overrides 'showstringspaces'
  showstringspaces=false,          % underline spaces within strings only
  showtabs=false,                  % show tabs within strings adding particular underscores
  stringstyle=\color{red!40},         % string literal style
  tabsize=2,                       % sets default tabsize to 2 spaces
  moredelim=**[is][\only<+>{\color{black}\lstset{style=highlight}}]{@}{@},
  moredelim=**[is][\only<1->{\color{black}\lstset{style=highlight}}]{@1}{@},
  moredelim=**[is][\only<2->{\color{black}\lstset{style=highlight}}]{@2}{@},
  moredelim=**[is][\only<3->{\color{black}\lstset{style=highlight}}]{@3}{@},
  moredelim=**[is][\only<4->{\color{black}\lstset{style=highlight}}]{@4}{@},
  moredelim=**[is][\only<5->{\color{black}\lstset{style=highlight}}]{@5}{@},
  moredelim=**[is][\only<6->{\color{black}\lstset{style=highlight}}]{@6}{@},
  moredelim=**[is][\only<7->{\color{black}\lstset{style=highlight}}]{@7}{@},
  moredelim=**[is][\only<8->{\color{black}\lstset{style=highlight}}]{@8}{@},
  moredelim=**[is][\only<9->{\color{black}\lstset{style=highlight}}]{@9}{@},
}

% Listening com destaque para comentarios (NAO ESTA FUNCIONANDO, TENTATIVA NA AULA 04
\newenvironment{lstlistingCMT}
  {\begin{lstlisting}[commentstyle=\color{red}\bfseries]}
  {\end{lstlisting}}


% Modificando espaçamento entre itens
%\newenvironment{myitemize}
%  {\begin{itemize}\setlength{\itemsep}{0.5\baselineskip}}
%  {\end{itemize}}
%
%\newenvironment{myenumerate}
%  {\begin{enumerate}\setlength{\itemsep}{0.5\baselineskip}}
%  {\end{enumerate}}

\newlength{\wideitemsep}
\setlength{\wideitemsep}{0.5\itemsep}
\addtolength{\wideitemsep}{6pt}
\let\olditem\item
\renewcommand{\item}{\setlength{\itemsep}{\wideitemsep}\olditem}



% Setup links
\definecolor{links}{HTML}{2A1B81}
\hypersetup{colorlinks,linkcolor=,urlcolor=links}


%Equacao bold em section
\newcommand{\eq}[1]{\texorpdfstring{$#1$}{}}

%Equacao bold fora de section
\usepackage{bm}
\newcommand{\eqbf}[1]{$\bm{#1}$} %\(\mathbf{#1}\)

%Alert bold
\newcommand{\alertbf}[1]{\alert{\textbf{#1}}}

%Inserir link para animação
\newcommand{\animation}[1]{
  \vskip5mm
  \begin{center}
    \href{#1}{\includegraphics[scale=0.1]{setup/animacao.jpg} \\ \footnotesize{Animação}}
  \end{center}
}

%Inserir link para vídeo
\newcommand{\video}[1]{
  \vskip5mm
  \begin{center}
    \href{#1}{\includegraphics[scale=0.1]{setup/animacao.jpg} \\ \footnotesize{Vídeo}}
  \end{center}
}

%Inserir link para animação e vídeo
\newcommand{\animationvideo}[2]{
  %\vskip5mm
  \begin{center}
    \begin{table}[h]
      \begin{tabular}{c c}
        \href{#1}{\includegraphics[scale=0.1]{setup/animacao}} & \href{#2}{\includegraphics[scale=0.1]{setup/animacao}} \\
        \href{#1}{\footnotesize{Animação}}                  & \href{#2}{\footnotesize{Vídeo}}
      \end{tabular}
    \end{table}
  \end{center}
}


% Tabelas
\newcommand{\mr}[2]{\multirow{#1}{*}{#2}}
\newcommand{\mc}[2]{\multicolumn{#1}{|c|}{#2}}



\institute[UFOP]{
%  \inst{1}%
\title[Autor]{ASN}

  Universidade Federal de Ouro Preto, UFOP
  \vskip0mm
  Departamento de Computação, DECOM
%  \vskip3mm
%   Website: \textbf{www.decom.ufop.br/anascimento}
  \vskip0mm
    Professor : \textbf{Anderson Almeida Ferreira}
  }

% Author, Title, etc.
\author[psilva]{\textbf{Felipe Braz, Matheus Peixoto, Pedro Henrique Oliveira}}
\title[BCC241]{BCC241 - Projeto e Análise de Algoritmos}
\date{
% \the\year{}
\vskip4mm
\includegraphics[scale=0.5]{setup/logos.jpg}
}

% Para inserir um frame de conteúdo antes de cada seção
\AtBeginSection[]{
    % \begin{frame}\frametitle{Conteúdo}
    %     \begin{small}
    %         \tableofcontents[currentsection]%, hideothersubsections] % PROCURAR SOLUCAO PARA APRESENTAR SUBSECOES
    %     \end{small}
    % \end{frame}
    \begin{frame}
        \vfill
        \centering
            \begin{beamercolorbox}[sep=8pt,center,shadow=true,rounded=true]{title}
            \usebeamerfont{title}\insertsectionhead\par%
        \end{beamercolorbox}
        \vfill
    \end{frame}
}

\newcommand{\begindocument}[1]{
  {\setbeamertemplate{footline}{} % Para não colocar rodapé na primeira página
  \begin{frame}[plain]%para não colocar barra de navegação na primeira página
    \titlepage
  \end{frame}
  }
  #1
  \begin{frame}{Conteúdo}
    \begin{small}
        \tableofcontents
    \end{small}
  \end{frame}
}
